\section{Introduction}\label{sec:intro}

Large language models predict next tokens; they do not enforce
programming-language rules. Even after billions of training tokens of
high-quality code, instruction-tuned code LMs hallucinate
ill-typed expressions, undefined identifiers, ill-formed borrows, and
nonsensical control flow. The hallucinations are not bugs in
particular models --- they are a structural consequence of how
autoregressive sampling commits to one token at a time without
look-back: any path the model places non-trivial probability mass on
becomes reachable, and a single bad token can cascade
\citep{jiang2024rocode,wang2025semguard}. Recent work on the
\emph{intrinsic} sources of hallucination in language modelling
confirms this: a frozen LM that has been trained on data with even a
small fraction of subtly incorrect programs will reliably generate
some subtle errors at inference, because next-token cross-entropy
gives no mechanism to refuse paths that violate global constraints
\citep{olausson2024selfrepair}.

Programming-language verifiers are exactly the missing mechanism.
A compiler, a type checker, a borrow checker, or a Language Server
Protocol (LSP) server, taken on the decidable fragment of code it
covers --- syntax, type, name resolution, lifetimes --- returns a
ground-truth verdict: either the program is well-formed or it is not.
This verdict is decoupled from how the program was produced. It is
also exactly the kind of side information a generative LM has no
internal way to consume. The natural design move is to put one in the
decoding loop.

Several recent systems do exactly this. Monitor-Guided Decoding
(MGD; \citealp{agrawal2023monitor}) consults an LSP at hand-picked
syntactic triggers --- the \code{.} that begins a member-access
expression in Java --- and masks the LM's next-token logits to those
prefixing a type-consistent identifier. ROCODE
\citep{jiang2024rocode} runs the compiler incrementally after each
generated statement, rolls back the buffer on an error, and decodes
the retry under an exponentially-decaying token-penalty mask.
SemGuard \citep{wang2025semguard} swaps the compiler for a learned
1.3B semantic classifier that scores each emitted line and triggers
line-level rollback below a threshold. Type-Constrained Code
Generation \citep{mundler2025typeconstrained} formalizes per-token
soundness in TypeScript by composing a prefix automaton with a
type-reachability search.

All four share one design choice: the verifier runs
\emph{synchronously} on the decoding critical path. The LM stops; the
verifier returns; the LM resumes. This is invisible when the
verifier is fast --- Python's \code{compile()} parses a partial
program in $\sim$0.1\,ms, $g{+}{+}\;\code{-fsyntax-only}$ parses a
warm translation unit in ${\sim}3\,\mathrm{ms}$ (though its
per-invocation subprocess cost inside a decoding loop is far
higher --- $\sim$$0.6$\,s amortized in our setup,
\S\ref{sec:headline-contrast}), and a 1.3B classifier evaluates one
line of context in single-digit milliseconds on a modern GPU. ROCODE's
benchmarks are Python and C++ \citep{jiang2024rocode}, SemGuard's are
Python and Java \citep{wang2025semguard}, MGD's are Java with
hand-picked dereference triggers \citep{agrawal2023monitor}. In each
case the per-call verifier cost is well below the LM's own per-token
latency, so the slowdown is tolerable. MGD reports a $\sim$83\%
mean-decoding-time slowdown \citep{agrawal2023monitor};
Type-Constrained reports a $39$--$52\%$ median per-instance
overhead \citep{mundler2025typeconstrained}. The cost is real but
amortized away by the gain in code quality.

The same design fails to scale to slower oracles. \code{cargo check}
on a typical Rust function takes $\sim$$35$--$500$
milliseconds incrementally (amortizing to $\sim$$0.18$\,s per call
in our setup); \code{javac} on a single Java method takes
${\sim}180\,\mathrm{ms}$ once JVM startup is amortized; a
whole-crate borrow check on an unfamiliar trait expansion can take
seconds. A synchronous decoding loop that blocks on such an oracle
spends most of its wall time idle. The two natural responses --- (1)
substitute a faster, weaker oracle, e.g.\ a learned classifier
\citep{wang2025semguard}; or (2) accept a $>10\times$ slowdown and
amortize against the larger code-quality gain --- both surrender
something the design originally promised. The first sacrifices
soundness (a learned classifier is wrong on a defined fraction of
inputs by construction); the second sacrifices the wall-clock budget
the verifier-guided pipeline was meant to outperform.

This paper takes a different response: \emph{run the verifier
asynchronously}. We introduce \textbf{SoundCode}, a producer-consumer
decoding loop in which:
\begin{itemize}[itemsep=2pt,topsep=2pt]
\item the LM streams tokens uninterrupted (the producer);
\item at each depth-zero structural boundary --- \code{;} or \code{\}}
  at brace-depth zero in C-family languages, newline at indent zero in
  Python --- a background task is spawned to run the verifier on the
  current buffer snapshot;
\item a consumer task drains finished checks in FIFO order;
\item on the first \emph{blocking} diagnostic, pending checks are
  cancelled, the buffer is truncated to the most recent
  verifier-blessed checkpoint, and the LM is re-prompted from there.
\end{itemize}
Crucially, the same rollback algorithm underlies the synchronous
ROCODE-style baseline that we compare against; \emph{the only
architectural difference is whether the verifier blocks the
decoder.} This isolates the sync-vs-async axis at fixed verifier,
fixed model, and fixed benchmark.

We evaluate SoundCode against ROCODE-style synchronous verification
on HumanEval-rs ($n\!=\!156$ after filtering) and HumanEval-cpp
($n\!=\!161$) from MultiPL-E \citep{cassano2022multipl}, using
Qwen2.5-Coder-7B-Instruct \citep{hui2024qwen25coder} via vLLM. We
report pass@1, mean wall-clock time per problem, mean rollback
count, and compile rate, plus an ablation on rollback target
(naive last-checkpoint vs.\ ROCODE's max-entropy formula).

\paragraph{Headline result.}
Holding the model, the verifier, and the rollback algorithm fixed,
SoundCode is $1.59\times$ faster on Rust ($1.30$\,s sync $\to$
$0.82$\,s async; $1.48\times$ at the per-problem median; async beats
sync on $116$ of $156$ paired problems, $74.4\%$) and $2.95\times$
faster on C++ ($3.65$\,s $\to$ $1.23$\,s; $2.84\times$ at the median;
async wins on $159$ of $161$ paired problems, $98.8\%$) than the
synchronous ROCODE-style equivalent. pass@1, by contrast, does not
separate the two schedulings: the per-language differences
($+5.0$\,pp for async on C++, $-3.2$\,pp on Rust) amount to a
handful of problems on a single greedy run, flip sign between
languages, and are not statistically significant (McNemar exact
$p\!=\!0.13$ and $p\!=\!0.30$ respectively). The claim of this paper
is therefore a \emph{latency} claim, not a quality claim:
asynchronous scheduling removes the synchronous verifier's
wall-clock penalty at statistically unchanged accuracy.

\paragraph{Contributions.} We frame the contribution along four axes:
\begin{enumerate}[itemsep=2pt,topsep=2pt]
\item \textbf{The first asynchronous design for verifier-guided code
  generation.} Every prior boundary-rollback system
  \citep{jiang2024rocode,wang2025semguard,ugare2025itergen} runs its
  verifier synchronously; every prior asynchronous-parallel-verify
  system \citep{leviathan2023speculative,chen2023speculative,fu2024lookahead}
  uses the LM itself as the verifier and gives up soundness.
  SoundCode occupies the open ``semantic-verifier $\times$
  structural-rollback $\times$ async'' design point.
\item \textbf{A controlled head-to-head measurement.} At fixed model,
  fixed verifier, and fixed rollback algorithm --- so the only axis
  that varies is whether the verifier blocks the decoder --- async
  delivers a $1.59\times$ wall-time speedup on Rust ($116/156$
  paired wins) and a $2.95\times$ speedup on C++ ($159/161$ paired
  wins), with pass@1 statistically indistinguishable between the
  schedulings on both languages.
\item \textbf{A verifier-cost-spectrum framing that predicts when
  async helps.} We sort verifiers along a cost spectrum (Python's
  \code{compile()} $\sim$$0.1\,\mathrm{ms}$; $g{+}{+}\;
  \code{-fsyntax-only}$ $\sim$$3\,\mathrm{ms}$ warm, amortizing to
  $\sim$$0.6$\,s per in-loop invocation in our setup;
  \code{cargo check}
  $\sim$$35$--$500\,\mathrm{ms}$, amortizing to
  $\sim$$0.18$\,s; \code{javac}
  $\sim$$180\,\mathrm{ms}$;
  whole-crate analyses up to seconds) and argue that the async design
  pays off whenever per-call verifier cost exceeds per-token decoding
  latency. The Rust-vs-C++ asymmetry in our results --- a smaller
  win on Rust, whose verifier amortizes cheaper in-loop
  ($\sim$$0.18$\,s/call vs C++'s $\sim$$0.6$\,s/call) --- is the
  first empirical confirmation of this prediction.
\item \textbf{A reproducible open-source implementation and live
  demo.} The producer is vLLM's \code{AsyncLLMEngine}
  \citep{vllm2023} run with \code{enforce\_eager=False}; the
  verifier is \code{cargo check} /
  $g{+}{+}\;\code{-fsyntax-only}$ via subprocess (the driver used
  for every measurement in this paper). An LSP driver
  (rust-analyzer via \code{multilspy}
  \citep{agrawal2023monitor,multilspy}) is implemented but not
  evaluated here (\S\ref{sec:setup}). Engineering effort
  to port to a new language is one boundary-detector and one
  verifier-driver class. Source:
  \href{https://github.com/leobwang/soundcode}%
       {\code{github.com/leobwang/soundcode}};
  live demo:
  \href{https://soundcode-demo.psixyzt.com}%
       {\code{soundcode-demo.psixyzt.com}}
  \citep{soundcode2025github}.
\end{enumerate}

The Olausson et al.\ matched-budget framework
\citep{olausson2024selfrepair} sets the bar any verifier-guided
system must clear: a self-repair loop only beats matched-budget
i.i.d.\ resampling when the feedback signal is \emph{strictly
stronger} than the generator. We argue that a real compiler is
exactly that. A learned classifier is not. SemGuard
\citep{wang2025semguard} acknowledges this trade-off explicitly. Our
async design lets us pay for the strictly-stronger oracle in
wall-clock terms without paying for it in throughput.
